\corrigesChapitre{4}{Dérivation et optimisation}
\begin{corrige}{1}{Règles usuelles}
\begin{questions}
\item $12x^3-2$.
\item $2(2x-1)\times2=8x-4$.
\item $3\mathrm e^{3x+2}$, par composition.
\end{questions}
\end{corrige}
\begin{corrige}{2}{Quotient et produit}
\begin{questions}
\item $\mathrm e^x+x\mathrm e^x=(x+1)\mathrm e^x$.
\item Pour $x\ne-2$, la dérivée vaut $[(x+2)-(x+1)]/(x+2)^2=1/(x+2)^2$.
\item $2x/(1+x^2)$ sur $\R$.
\end{questions}
\end{corrige}
\begin{corrige}{3}{Tangente}
\begin{questions}
\item $f(2)=5$ et $f\,\!^{\prime}(2)=4$.
\item $y=5+4(x-2)=4x-3$.
\item $f(2{,}01)\simeq5{,}04$ ; la valeur exacte est $5{,}0401$, soit une erreur de $0{,}0001$.
\end{questions}
\end{corrige}
\begin{corrige}{4}{Signe et tableau}
\begin{questions}
\item $f\,\!^{\prime}(x)=2x-6$, nulle en $3$.
\item Décroissance sur $[0,3]$, puis croissance sur $[3,5]$.
\item Minimum $f(3)=-4$ ; aux bornes $f(0)=5$, $f(5)=0$, donc maximum $5$.
\end{questions}
\end{corrige}
\begin{corrige}{5}{Rectangle}
\begin{questions}
\item L’autre longueur vaut $30-x$, donc $A=x(30-x)$, avec $0<x<30$ pour un rectangle non dégénéré.
\item $A\,\!^{\prime}=30-2x$ change de $+$ à $-$ en $15$.
\item Le carré de côté $15$ m maximise l’aire, qui vaut $225$ m$^2$.
\end{questions}
\end{corrige}
\begin{corrige}{6}{Bénéfice}
\begin{questions}
\item $B(x)=-x^2+30x-100$.
\item $B\,\!^{\prime}=30-2x$ : maximum en $x=15$, égal à $125$.
\item Les racines sont $15\pm5\sqrt5$ ; le bénéfice est positif strictement entre ces deux racines, soit environ $3{,}82<x<26{,}18$.
\end{questions}
\end{corrige}
\begin{corrige}{7}{Coût moyen}
\begin{questions}
\item $C_m(x)=2x+18+200/x$.
\item $C_m\,\!^{\prime}=2-200/x^2$ change de signe en $10$ ; minimum $58$.
\item $C\,\!^{\prime}(10)=58$ ; la différence exacte vaut $2(121-100)+18=60$. Le coût marginal est ici une approximation.
\end{questions}
\end{corrige}
\begin{corrige}{8}{Une exponentielle composée}
\begin{questions}
\item $f\,\!^{\prime}(x)=(1-x)\mathrm e^{-x}$.
\item L’exponentielle étant positive, $f$ croît jusqu’à $1$ puis décroît.
\item Le maximum vaut $f(1)=\mathrm e^{-1}$, atteint en $x=1$.
\end{questions}
\end{corrige}
\begin{corrige}{9}{Une boîte sans couvercle}
\begin{questions}
\item $V(x)=x(12-2x)^2$ pour $0<x<6$.
\item $V\,\!^{\prime}(x)=(12-2x)(12-6x)$.
\item Sur le domaine, le premier facteur est positif ; le signe change en $x=2$. Le maximum vaut $2\times8^2=128$ cm$^3$.
\end{questions}
\end{corrige}
\begin{corrige}{10}{Dérivée nulle : attention}
\begin{questions}
\item $3x^2\geq0$ des deux côtés ; elle s’annule seulement en zéro.
\item Non : les valeurs sont négatives à gauche, positives à droite et la fonction reste croissante.
\item La dérivée $4x^3$ change de $-$ à $+$ : zéro est un minimum.
\end{questions}
\end{corrige}
\begin{corrige}{11}{Production entière}
\begin{questions}
\item $B\,\!^{\prime}=-4x+19$ s’annule en $19/4=4{,}75$.
\item $B(4)=24$, $B(5)=25$.
\item La fonction croît avant $4{,}75$ puis décroît : l’entier optimal est $5$ avec un bénéfice de $25$ unités monétaires.
\end{questions}
\end{corrige}
\begin{corrige}{12}{Deux disciplines, un modèle}
\begin{questions}
\item $v(t)=s\,\!^{\prime}(t)=20-4t$, en m/s.
\item La vitesse est positive avant $5$ puis nulle à $5$ ; le maximum sur $[0,5]$ vaut $s(5)=50$ m.
\item Le domaine de validité est imposé par le phénomène. Le prolongement prédirait ensuite un mouvement de retour qui doit être justifié physiquement.
\end{questions}
\end{corrige}
